\documentclass[11pt,a4paper]{article}
\usepackage[utf8]{inputenc}
\usepackage[vietnamese]{babel}
\usepackage{amsmath,amssymb}
\usepackage{geometry}
\geometry{a4paper, total={180mm,260mm}, left=15mm, right=15mm, top=18mm, bottom=18mm}
\usepackage{tikz}
\usetikzlibrary{shapes.geometric, arrows.meta, positioning, calc}
\usepackage{tcolorbox}
\tcbuselibrary{skins,breakable}
\usepackage{xcolor}
\usepackage{multicol}
\usepackage{array}
\usepackage{fontawesome}
\usepackage{fancyhdr}

\pagestyle{fancy}
\fancyhf{}
\renewcommand{\headrulewidth}{0.5pt}
\renewcommand{\footrulewidth}{0.5pt}
\fancyhead[L]{\small\bfseries\color{teal}\faBook\ SỔ TAY KIẾN THỨC TOÁN 3}
\fancyhead[R]{\small\bfseries\color{teal}Lớp Toán Thầy Lĩnh -- loptoanthaylinh.vn}
\fancyfoot[C]{\small\color{gray}-- Sổ tay kiến thức: Bài 41. Ôn tập phép nhân, phép chia trong phạm vi 100, 1 000 --}

\definecolor{primary}{HTML}{005F73}
\definecolor{secondary}{HTML}{0A9396}
\definecolor{accent}{HTML}{EE9B00}
\definecolor{danger}{HTML}{AE2012}
\definecolor{boxbg}{HTML}{F4F6F8}

\begin{document}

% --- BANNER TIÊU ĐỀ ---
\begin{tcolorbox}[
  enhanced,
  colback=white,
  colframe=primary,
  arc=3mm,
  boxrule=1.2pt,
  left=4mm, right=4mm, top=2.5mm, bottom=2.5mm,
  drop shadow=black!10!white
]
  \centering
  {\Large\bfseries\color{primary}\faBookmark\ BÀI 41: ÔN TẬP PHÉP NHÂN, CHIA TRONG PHẠM VI 100, 1 000}\\[1.5mm]
  {\small\itshape\color{black!75}Ôn tập học kì 1: Toàn bộ bảng nhân chia 2 - 9, nhân chia số có 2, 3 chữ số trên 1 trang A4}
\end{tcolorbox}

\vspace{-2mm}


\noindent
\begin{minipage}[t]{0.49\textwidth}\vspace{0pt}
\begin{tcolorbox}[
  enhanced,
  colback=blue!4!white,
  colframe=blue!75!black,
  arc=2mm, boxrule=0.8pt,
  left=2mm, right=2mm, top=1.8mm, bottom=1.8mm,
  title={\footnotesize\bfseries\color{white}\faKey\ 1. Bảng nhân chia và tính chất}
]
{\footnotesize
\begin{itemize}\setlength{\itemsep}{1.5pt}\setlength{\leftskip}{-2.5mm}
  \item \textbf{Thuộc lòng 8 bảng cửu chương:} Bảng nhân, chia từ $2$ đến $9$.
  \item \textbf{Tính chất đặc biệt:}
  \begin{itemize}\setlength{\itemsep}{0.5pt}\setlength{\leftskip}{-2.5mm}
    \item $a \times 1 = a$; $a : 1 = a$.
    \item $a \times 0 = 0$; $0 : a = 0\ (a \ne 0)$.
    \item Phép nhân có tính chất giao hoán: $a \times b = b \times a$.
  \end{itemize}
  \item \textbf{Tìm thành phần chưa biết:}
  $$x \times a = b \to x = b : a \qquad a : x = b \to x = a : b$$
\end{itemize}
}
\end{tcolorbox}
\end{minipage}
\hfill
\begin{minipage}[t]{0.49\textwidth}\vspace{0pt}
\begin{tcolorbox}[
  enhanced,
  colback=teal!4!white,
  colframe=teal!75!black,
  arc=2mm, boxrule=0.8pt,
  left=2mm, right=2mm, top=1.8mm, bottom=1.8mm,
  title={\footnotesize\bfseries\color{white}\faCalculator\ 2. Kĩ thuật đặt tính phép chia}
]
{\footnotesize
\begin{itemize}\setlength{\itemsep}{1.5pt}\setlength{\leftskip}{-2.5mm}
  \item \textbf{Phép chia hết:} Số dư bằng $0$ ($r = 0$).
  \item \textbf{Phép chia có dư:}
  $$\text{Số bị chia} = \text{Thương} \times \text{Số chia} + r \quad (r < d)$$
  \item Số dư lớn nhất luôn bằng: $\text{Số chia} - 1$.
  \item Nhân số tròn chục, tròn trăm: nhân phần số rồi thêm các chữ số $0$ vào tận cùng bên phải tích.
\end{itemize}
}
\end{tcolorbox}
\end{minipage}

\vspace{-2.5mm}
% --- KHỐI 3: SƠ ĐỒ TƯ DUY TIKZ ---
\noindent
\begin{tcolorbox}[
  enhanced,
  colback=orange!4!white,
  colframe=orange!80!black,
  arc=2mm, boxrule=0.8pt,
  left=2mm, right=2mm, top=1.5mm, bottom=1.5mm,
  title={\footnotesize\bfseries\color{white}\faSitemap\ SƠ ĐỒ TƯ DUY: CÂY KẾT NỐI KIẾN THỨC NHÂN CHIA HK1}
]
\centering
\begin{tikzpicture}[
  box/.style={draw, rounded corners=2pt, align=center, font=\scriptsize\bfseries, inner sep=2.5pt, minimum height=0.75cm},
  arr/.style={-{Stealth[length=2mm]}, thick}
]
  \node[box, fill=purple!15, draw=purple!80!black, text width=2.6cm] (k1) at (0, 0) {BẢNG NHÂN CHIA\\Bảng 2 đến 9\\Mối quan hệ chéo};
  \node[box, fill=teal!15, draw=teal!80!black, text width=2.6cm] (k2) at (3.5, 0) {ĐẶT TÍNH\\$2, 3$ chữ số $\times / : 1$\\Nhân phải, Chia trái};
  \node[box, fill=blue!12, draw=blue!80!black, text width=2.6cm] (k3) at (7.0, 0) {CHIA CÓ DƯ\\Số dư $r < \text{số chia}$\\$r$ lớn nhất $= d - 1$};
  \node[box, fill=orange!15, draw=orange!80!black, text width=2.6cm] (k4) at (10.5, 0) {TOÁN LỜI VĂN\\Gấp / Giảm số lần\\Toán 2 bước tính};

  \draw[arr, purple!80!black] (k1) -- (k2);
  \draw[arr, teal!80!black] (k2) -- (k3);
  \draw[arr, blue!80!black] (k3) -- (k4);
\end{tikzpicture}
\end{tcolorbox}

\vspace{-2.5mm}
% --- KHỐI 4: VÍ DỤ & MẸO TRÁNH BẪY ---
\noindent
\begin{minipage}[t]{0.49\textwidth}\vspace{0pt}
\begin{tcolorbox}[
  enhanced,
  colback=green!4!white,
  colframe=green!60!black,
  arc=2mm, boxrule=0.8pt,
  left=2mm, right=2mm, top=1.8mm, bottom=1.8mm,
  title={\footnotesize\bfseries\color{white}\faCheckCircle\ Ví dụ thực hành chuẩn mực}
]
{\footnotesize
\textbf{1. Tính nhẩm nhanh:}\\
$300 \times 3 = \mathbf{900}$.\\
$800 : 4 = \mathbf{200}$.\\
$40 \times 5 = \mathbf{200}$.\\
\textbf{2. Tìm số bị chia khi biết số dư:}\\
Tìm $x$ biết $x : 7 = 14$ (dư lớn nhất).\\
Số dư lớn nhất là: $7 - 1 = 6$.\\
Số bị chia: $x = 14 \times 7 + 6 = 98 + 6 = \mathbf{104}$.
}
\end{tcolorbox}
\end{minipage}
\hfill
\begin{minipage}[t]{0.49\textwidth}\vspace{0pt}
\begin{tcolorbox}[
  enhanced,
  colback=red!4!white,
  colframe=red!70!black,
  arc=2mm, boxrule=0.8pt,
  left=2mm, right=2mm, top=1.8mm, bottom=1.8mm,
  title={\footnotesize\bfseries\color{white}\faExclamationTriangle\ Mẹo tránh bẫy trắc nghiệm}
]
{\footnotesize
\begin{itemize}\setlength{\itemsep}{1.5pt}\setlength{\leftskip}{-2.5mm}
  \item \textbf{Bẫy quên cộng số dư khi tìm số bị chia:}\\
  Rất nhiều bạn lấy $14 \times 7 = 98$ rồi dừng lại, quên cộng số dư $6$! Công thức chuẩn: $\text{Số bị chia} = \text{Thương} \times \text{Số chia} + \text{Dư}$.
  \item Không bao giờ có phép chia cho số $0$!
\end{itemize}
}
\end{tcolorbox}
\end{minipage}


\end{document}
